% GENERATED FILE. Edit canonical lessons, then run npm run generate:books.
% canonical-source-hash: fnv1a64:7d80770b7027694a
% canonical-lessons: FR-C458-dune, FR-C458-maree, FR-C458-crique, FR-C458-recif, FR-C458-coquillage

\chapter{Au bord de la mer}
\label{ch:fr-au-bord-de-la-mer}
\hlchaptermodality{458}

\begin{chapteropening}
\textbf{By the end of this chapter:} \emph{I can describe the coast: the dune, the tide, the cove, the reef and the seashell.}

Complete the last lesson of chapter 458: I can describe the coast: the dune, the tide, the cove, the reef and the seashell.
\end{chapteropening}

\section[la dune]{la dune — the dune}

\label{lesson:FR-C458-dune}



\begin{quote}
Before the new one: say the French for the moss, then the French for the bark.
\end{quote}

\subsection*{You'll want to know: la dune}
\textbf{la dune} — \textquotedblleft{}the dune\textquotedblright{}.

\begin{quote}
Derrière la plage, il y a de hautes dunes.
\end{quote}

\begin{quote}
\emph{Behind the beach there are high dunes.}
\end{quote}

\emph{Une dune de sable} — a sand dune.

\begin{cousinweb}
\emph{dune} was borrowed from Middle Dutch \emph{dūne} (EtymDB).
\end{cousinweb}

\begin{grammarlens}[title={What you've built}]
One more word for telling, explaining and describing at B1.
\end{grammarlens}

\subsection*{Your turn}
\begin{itemize}
\raggedright
  \item \emph{Say it:} \emph{la dune}
  \item \emph{Say it:} \emph{la dune}, once more
  \item \emph{Say it:} a sentence of your own with \emph{la dune}
  \item \emph{Recall:} say the French for the moss, then the French for the bark, then say \emph{la dune} again
\end{itemize}

\begin{tcolorbox}[breakable,colback=teal!4,colframe=teal!35!black,title={Before you move on}]
What does la dune mean? (\textquotedblleft{}The dune\textquotedblright{}.) Say it once more.
\end{tcolorbox}
\section[la marée]{la marée — the tide}

\label{lesson:FR-C458-maree}



\begin{quote}
Before the new one: say the French for the bark, then the French for the dune.
\end{quote}

\subsection*{You'll want to know: la marée}
\textbf{la marée} — \textquotedblleft{}the tide\textquotedblright{}.

\begin{quote}
À marée basse, on peut marcher jusqu'à l'île.
\end{quote}

\begin{quote}
\emph{At low tide you can walk to the island.}
\end{quote}

\emph{Marée haute} — high tide; \emph{une marée noire} — an oil spill.

\begin{cousinweb}
\emph{marée} is built on \emph{mer} 'sea' (EtymDB), which is inherited from Latin \emph{mare}.
\end{cousinweb}

\begin{grammarlens}[title={What you've built}]
One more word for telling, explaining and describing at B1.
\end{grammarlens}

\subsection*{Your turn}
\begin{itemize}
\raggedright
  \item \emph{Say it:} \emph{la marée}
  \item \emph{Say it:} \emph{la marée}, once more
  \item \emph{Say it:} a sentence of your own with \emph{la marée}
  \item \emph{Recall:} say the French for the bark, then the French for the dune, then say \emph{la marée} again
\end{itemize}

\begin{tcolorbox}[breakable,colback=teal!4,colframe=teal!35!black,title={Before you move on}]
What does la marée mean? (\textquotedblleft{}The tide\textquotedblright{}.) Say it once more.
\end{tcolorbox}
\section[la crique]{la crique — the cove, the creek}

\label{lesson:FR-C458-crique}



\begin{quote}
Before the new one: say the French for the dune, then the French for the tide.
\end{quote}

\subsection*{You'll want to know: la crique}
\textbf{la crique} — \textquotedblleft{}the cove, the creek\textquotedblright{}.

\begin{quote}
Nous avons trouvé une petite crique sans personne.
\end{quote}

\begin{quote}
\emph{We found a little cove with nobody in it.}
\end{quote}

A small, sheltered beach between rocks.

\begin{cousinweb}
\emph{crique} was borrowed from Old Norse \emph{kriki} 'bend, nook' (EtymDB).
\end{cousinweb}

\begin{grammarlens}[title={What you've built}]
One more word for telling, explaining and describing at B1.
\end{grammarlens}

\subsection*{Your turn}
\begin{itemize}
\raggedright
  \item \emph{Say it:} \emph{la crique}
  \item \emph{Say it:} \emph{la crique}, once more
  \item \emph{Say it:} a sentence of your own with \emph{la crique}
  \item \emph{Recall:} say the French for the dune, then the French for the tide, then say \emph{la crique} again
\end{itemize}

\begin{tcolorbox}[breakable,colback=teal!4,colframe=teal!35!black,title={Before you move on}]
What does la crique mean? (\textquotedblleft{}The cove, the creek\textquotedblright{}.) Say it once more.
\end{tcolorbox}
\section[le récif]{le récif — the reef}

\label{lesson:FR-C458-recif}



\begin{quote}
Before the new one: say the French for the tide, then the French for the cove.
\end{quote}

\subsection*{You'll want to know: le récif}
\textbf{le récif} — \textquotedblleft{}the reef\textquotedblright{}.

\begin{quote}
Les bateaux évitent le récif près du port.
\end{quote}

\begin{quote}
\emph{The boats steer clear of the reef near the harbour.}
\end{quote}

\emph{Un récif de corail} — a coral reef.

\begin{cousinweb}
\emph{récif} was borrowed from Spanish \emph{arrecife}, which is usually traced to Arabic \emph{ar-raṣīf} 'causeway, raised way'. EtymDB lists it without a source.
\end{cousinweb}

\begin{grammarlens}[title={What you've built}]
One more word for telling, explaining and describing at B1.
\end{grammarlens}

\subsection*{Your turn}
\begin{itemize}
\raggedright
  \item \emph{Say it:} \emph{le récif}
  \item \emph{Say it:} \emph{le récif}, once more
  \item \emph{Say it:} a sentence of your own with \emph{le récif}
  \item \emph{Recall:} say the French for the tide, then the French for the cove, then say \emph{le récif} again
\end{itemize}

\begin{tcolorbox}[breakable,colback=teal!4,colframe=teal!35!black,title={Before you move on}]
What does le récif mean? (\textquotedblleft{}The reef\textquotedblright{}.) Say it once more.
\end{tcolorbox}
\section[le coquillage]{le coquillage — the seashell; the shellfish}

\label{lesson:FR-C458-coquillage}



\begin{quote}
Before the new one: say the French for the cove, then the French for the reef.
\end{quote}

\subsection*{You'll want to know: le coquillage}
\textbf{le coquillage} — \textquotedblleft{}the seashell; the shellfish\textquotedblright{}.

\begin{quote}
La petite fille ramasse des coquillages sur la plage.
\end{quote}

\begin{quote}
\emph{The little girl collects shells on the beach.}
\end{quote}

\emph{Les coquillages} are also shellfish to eat.

\begin{cousinweb}
\emph{coquillage} is built on \emph{coquille} 'shell' (EtymDB lists it), which goes back to Vulgar Latin \emph{conchilia}, from Latin \emph{concha} 'shellfish, shell'.
\end{cousinweb}

\begin{grammarlens}[title={What you've built}]
One more word for telling, explaining and describing at B1.
\end{grammarlens}

\subsection*{Your turn}
\begin{itemize}
\raggedright
  \item \emph{Say it:} \emph{le coquillage}
  \item \emph{Say it:} \emph{le coquillage}, once more
  \item \emph{Say it:} a sentence of your own with \emph{le coquillage}
  \item \emph{Recall:} say the French for the cove, then the French for the reef, then say \emph{le coquillage} again
\end{itemize}

\begin{tcolorbox}[breakable,colback=teal!4,colframe=teal!35!black,title={Before you move on}]
What does le coquillage mean? (\textquotedblleft{}The seashell; the shellfish\textquotedblright{}.) Say it once more.
\end{tcolorbox}
